\part{Going further}

\chapter{Scaling up: 350M to 1.3B}
\chapfig{fig_scaling.png}

Chapter~9 ended with a wall: the 350M model was capacity-bound on reasoning. The book's central prediction was
that \emph{scale} --- not more data or more training --- is the lever that moves it. So we tested the
prediction directly.

\section{Building a bigger model is two numbers}
We doubled the width (dimension 1024 $\rightarrow$ 2048), keeping 24 layers, giving \textbf{1.30 billion
parameters}. Same code, same data pipeline, same tokenizer. This is the payoff of having built everything
ourselves.

\section{Re-tuning for the new size}
A learning rate that is right for 350M is usually too high for 1.3B. Rather than guess, we ran a quick
three-model learning-rate ablation \emph{at the new scale} --- and indeed \texttt{1e-3} was now too hot;
\textbf{5e-4} won. This is the playbook applied to itself: a \$40 ablation before a \$880 run.

\section{The result}
\fig{plots/plot_scaling.pdf}{0.82}{The capacity lever, tested. At 350M, MedMCQA was stuck near chance no matter
what we did. At 1.3B it \emph{moved} --- the metric that data and training could not budge responded to scale,
exactly as predicted.}

\begin{center}\small
\begin{tabular}{@{}lccl@{}}
\toprule
\textbf{Metric} & \textbf{350M} & \textbf{1.3B} & \textbf{Read}\\
\midrule
MedMCQA (chance 0.25) & 0.290 & \textbf{0.321} & cleared chance --- scale moved it\\
PubMedQA & 0.608 & \textbf{0.654} & $+4.6$ points\\
PubMed val-loss $\downarrow$ & 2.27 & \textbf{2.12} & better\\
PMC val-loss $\downarrow$ & 2.13 & \textbf{1.97} & better\\
\bottomrule
\end{tabular}
\end{center}

\section{The generation jump is bigger than the table}
On \emph{``mechanism of action of metformin''}:
\begin{quote}\small
\textbf{350M:} ``Metformin is a \emph{glycoside antibiotic} similar to erythromycin\ldots'' --- fluent but
flatly wrong.\\[3pt]
\textbf{1.3B:} ``1. Insulin sensitization. 2. Decreased hepatic glucose production. 3. Decreased glucose
absorption\ldots'' --- a correct, structured mechanism.
\end{quote}

\begin{lesson}
Scale is the lever for reasoning. It moved the exam score that nothing else could --- but MedMCQA at 0.32 is
still modest: the wall \emph{lifts}, it does not vanish. For board-exam accuracy you keep climbing
(2.7B$+$). And the cheapest reliability win isn't scale at all --- it is retrieval (next chapter).
\end{lesson}

\begin{gotcha}[The bigger model ran out of memory]
The first 1.3B run crashed instantly with a CUDA out-of-memory error. The culprit was the logits tensor ---
\texttt{(batch, sequence, 32000)} --- which at sequence length 2048 is enormous, on top of the larger optimizer
state. The fix: a much smaller micro-batch with more gradient accumulation (same effective batch), plus
expandable memory segments. Big models need small micro-batches.
\end{gotcha}

\chapter{Retrieval: answers it can cite}
\chapfig{fig_rag.png}

A small model's worst habit is confident wrongness. The fix is not always a bigger model. Often it is
\textbf{retrieval}: instead of trusting the model's memory, fetch real source passages at question time and
have the model answer \emph{grounded in them}.

\section{What we built}
We indexed \textbf{300{,}000 PubMed passages} with a biomedical embedding model (PubMedBERT) into a FAISS
vector index. At query time we embed the question, retrieve the most similar passages, and prompt the 1.3B
model to answer using them --- returning the sources so a human can verify.

\fig{figures/fig_grounding.png}{0.55}{Every answer is now tethered to evidence. The model's response is
anchored to the retrieved passages, and those passages are shown with similarity scores.}

\section{It works --- with an honest caveat}
On \emph{``the role of AMPK in the antidiabetic action of metformin,''} retrieval pulled the right paper and
the model answered correctly and \emph{cited it}. But on a vaguer query (``mechanism of metformin'') retrieval
pulled tangential metformin-\emph{cancer} papers and the answer drifted. Retrieval quality matters.

\begin{numbers}[Why RAG is the highest-ROI next step]
It adds factual grounding and citations --- the thing enterprises actually need for trust --- for a few hundred
dollars, no model retraining, and it gets \emph{better} as the model gets better at reading. A 1.3B model
$+$ RAG beats a 7B model alone on factual reliability, at a fraction of the cost.
\end{numbers}

\begin{lesson}
When facts matter more than your budget can scale to, reach for retrieval before more parameters. Grounding
$+$ citations is the cheapest, highest-trust reliability layer for a domain assistant.
\end{lesson}
